\section{Global state, local instruments, and Bell's triadic correlations}
\label{chap:bell-local-global}

The question posed by Bell \cite{Bell1964} relates the local description of
outcomes to the global structure of their correlations. It asks whether a family of
joint probabilities admits a decomposition into local responses conditioned by the
same prior state. In HMT this question has a natural domain: the enriched state preserves
genealogy, orientation, carry, and boundary memory, while each
instrument publishes a local coordinate. The richness of the common state and the
probabilistic form of its readings constitute distinct mathematical data. The former
organizes the provenance of the correlation; the latter determines its position relative
to the local Bell polytope.

The purpose of this chapter is to traverse that complete chain. We shall begin with a
global preparation on the TPK state space, define the instruments of
the two wings, recover the CHSH bound, and prove its stability under
enrichment of the shared state; we shall then pass to the ternary alphabet intrinsic
to TRIT. In that new
domain the information unit is the trit and the pertinent inequality is CGLMP. The
internal elliptic structure further provides a real realization of the
order-three Fourier transform through the operator \(J_E^2=-I\). The result is
a precise formulation of the relation among holographic totality, local reading, and
quantum correlation.

\subsection{Enhanced preparation and boundary instruments}

Let \(X_{\rm enr}\) be the space of enriched TPK states, and let
\(\mu\) be a probability measure describing the common preparation. An element
\(x\in X_{\rm enr}\) comprises the local state together with its transport history.
The reading is introduced by means of instruments. For two
wings, Alice and Bob, with setting sets \(\mathcal A\) and \(\mathcal B\),
we write
\[
 \mathcal I_A^a:X_{\rm enr}\longrightarrow\mathcal D(\mathcal O_A),
 \qquad
 \mathcal I_B^b:X_{\rm enr}\longrightarrow\mathcal D(\mathcal O_B),
\]
where \(\mathcal D(\mathcal O)\) denotes the simplex of distributions on the
outcome alphabet \(\mathcal O\). For each prepared state, the kernels
\(\mathcal I_A^a(A\mid x)\) and \(\mathcal I_B^b(B\mid x)\) express the
probabilities published by the local readers.

The local factored form of the joint distribution is
\begin{equation}
 P_{\rm loc}(A,B\mid a,b)
 =\int_{X_{\rm enr}}
   \mathcal I_A^a(A\mid x)\,
   \mathcal I_B^b(B\mid x)\,\dd\mu(x).
 \label{eq:bell-hmt-factorizado}
\end{equation}
This equation preserves a common global state and allows that state to possess arbitrarily rich memory. Its characteristic assertion is the factorization of responses once \(x\) has been fixed. Independence from the settings further requires that the measure \(\mu\) be common to all choices of \(a\) and \(b\), and complete sampling requires the observed population to represent that same preparation for every pair of settings.

This distinction is linked directly to the HMT definition of measurement. The
reversible coupling between system and apparatus produces a joint state; the
instrument selects a partition of that state, and the reader publishes one of its
fibers. For a setting \(a\) and an outcome \(A\), the fiber
\[
 \Omega_{a,A}
 =\{x\in X_{\rm enr}:\operatorname{Out}_a(x)=A\}
\]
represents the set of histories compatible with the reading. Different settings
induce different partitions of the same space; the published random variable
is therefore typed by its instrumental context. Contextuality is thus formulated
in the same language of instruments, fibers, and survival.

\subsection{Local factorization, CHSH bound, and stability under enrichment}

First consider binary outcomes
\(\mathcal O_A=\mathcal O_B=\{-1,+1\}\) and two settings per wing. The
correlations and the Clauser--Horne--Shimony--Holt functional \cite{ClauserHorneShimonyHolt1969} are
\begin{equation}
 E_{ab}=\sum_{A,B=\pm1}AB\,P(A,B\mid a,b),
 \qquad
 S_{\rm CHSH}=E_{00}+E_{01}+E_{10}-E_{11}.
 \label{eq:chsh}
\end{equation}

\begin{theorem}[CHSH Bound under Local Factorization]
\label{thm:chsh-local}
If the joint distribution has the form
\eqref{eq:bell-hmt-factorizado}, the preparation is independent of the settings, and
sampling preserves the common population in the four pairs of settings, then
\[
 \lvert S_{\rm CHSH}\rvert\leq2.
\]
\end{theorem}

\begin{proof}
For each \(x\), define the conditional expectations
\[
 \overline A_a(x)=\sum_{A=\pm1}A\,\mathcal I_A^a(A\mid x),
 \qquad
 \overline B_b(x)=\sum_{B=\pm1}B\,\mathcal I_B^b(B\mid x).
\]
Ambas belong to \([-1,1]\). The integrando of
\(S_{\rm CHSH}\) is
\[
 \overline A_0(\overline B_0+\overline B_1)
 +\overline A_1(\overline B_0-\overline B_1).
\]
For any \(u,v\in[-1,1]\),
\(\lvert u+v\rvert+\lvert u-v\rvert\leq2\) holds. The absolute value of the
integrand therefore does not exceed \(2\). Integration with respect to a single
probability measure \(\mu\) preserves the bound.
\end{proof}

The theorem leads to a specifically important consequence for HMT.

\begin{corollary}[Persistence of the bound under state enrichment]
\label{cor:no-go-bell-memoria}
Replacing an elementary hidden variable with the complete state
\(x\in X_{\rm enr}\) preserves the bound
\(\lvert S_{\rm CHSH}\rvert\leq2\) provided that the two instruments continue to
factor as in \eqref{eq:bell-hmt-factorizado}.
\end{corollary}

\begin{proof}
The state \(x\) plays the role of the shared variable in the
\cref{thm:chsh-local}. The argument depends on factorization of the conditional
responses and on a preparation common to the four pairs of settings; the internal
dimension of \(X_{\rm enr}\) leaves both properties intact.
\end{proof}

This corollary positively establishes the locus at which an HMT theory of
quantum correlations acts. The global genealogy provides the common domain. To
realize a violation, the physical morphism must provide a joint
nonseparable rule. In this formulation, the Bell correlation expresses a property of the
joint representation of the state, compatible with the operational independence of
the marginals.

\subsection{Extension of binary entropy to the qutrit system and emergence of the invariant \(\log 3\)}

TRIT naturally selects the alphabet
\(\mathbb Z/3\mathbb Z\). For a distribution
\(p=(p_0,p_1,p_2)\), its entropy expressed in trits is
\begin{equation}
 H_3(p)=-\sum_{j=0}^{2}p_j\log_3p_j
       =\frac{H_e(p)}{\log 3}.
 \label{eq:entropia-base-tres}
\end{equation}
The same relation holds for mutual information; to avoid confusing it with the
Bell functional introduced shortly, we denote it by
\[
 \mathsf I^{(3)}(X:Y)
 =\frac{\mathsf I^{(e)}(X:Y)}{\log 3}.
\]

\begin{proposition}[Invariance of Locality Under Change of Information Unit]
\label{prop:log3-localidad}
The passage from natural logarithms to base-three logarithms multiplies entropies and
information by the positive constant \((\log 3)^{-1}\). It preserves the
joint distributions, their marginals, and the local factorization condition.
\end{proposition}

\begin{proof}
The identity \(\log_3u=(\log u)/(\log 3)\) proves the first assertion. The second follows
because change of basis acts on the value used to quantify an already given
distribution and leaves invariant both its terms and the stochastic kernels
that factorize it.
\end{proof}

The appearance of \(\log 3\) places the analysis in an information unit suited to the
ternary alphabet. The substantive change for Bell is the passage from two to three outcomes
per instrument. This passage leads from the binary CHSH functional to a
qutrit inequality and allows the TRIT structure to enter both the definition of the
settings and the choice of units.

\subsection{The CGLMP inequality for \(3\) outcomes}

Let \(A_1,A_2,B_1,B_2\) be variables taking values in
\(\mathbb Z/3\mathbb Z\); all equalities in this section are understood
modulo three. We define the
Collins--Gisin--Linden--Massar--Popescu functional \cite{CGLMP2002}
\begin{align}
 I_3^{\rm CGLMP}={}&P(A_1=B_1)+P(B_1=A_2+1)
 +P(A_2=B_2)+P(B_2=A_1)\notag\\
 &-P(A_1=B_1-1)-P(B_1=A_2)
 -P(A_2=B_2-1)-P(B_2=A_1-1).
 \label{eq:cglmp-tres}
\end{align}

\begin{theorem}[Local triadic bound]
\label{thm:cglmp-local}
Every distribution of four ternary instruments admitting a factorized local
model satisfies
\[
 I_3^{\rm CGLMP}\leq2.
\]
\end{theorem}

\begin{proof}
The set of local distributions is the convex hull of the deterministic
assignments, so it suffices to prove the bound at its vertices. For a fixed
assignment, write
\[
 x=A_1-B_1,\qquad y=B_1-A_2,\qquad
 z=A_2-B_2,\qquad t=B_2-A_1.
\]
The relation of closure is \(x+y+z+t=0\). In these terminos,
\[
\begin{aligned}
 I_3^{\rm CGLMP}={}&
 \mathbf1_{x=0}+\mathbf1_{y=1}
 +\mathbf1_{z=0}+\mathbf1_{t=0}\\
 &-\mathbf1_{x=2}-\mathbf1_{y=0}
 -\mathbf1_{z=2}-\mathbf1_{t=2}.
\end{aligned}
\]
The twenty-seven choices of \((x,y,z)\) determine \(t\). For
\(y=0\), the nine patterns produce \(-4\) once, \(-1\) seven times, and
\(2\) once; for \(y=1\), they produce \(-1\) three times and \(2\) six times;
for \(y=2\), they produce \(-1\) six times and \(2\) three times. The maximum at
every vertex is \(2\), and convexity preserves that bound.
\end{proof}

The preceding inequality describes the convex boundary of local models in the
same alphabet used by HMT. Its physical evaluation starts from a joint distribution obtained from a
preparation and four concrete instruments. The condition
\(I_3^{\rm CGLMP}>2\), when obtained with a complete sample and independently fixed
settings, excludes a local factorization of the form
\eqref{eq:bell-hmt-factorizado} for ternary outcomes.

\subsection{Triadic Transform from the Internal Elliptic Structure}

Realization of the qutrit settings requires a transformation that mixes the
three outputs. Let \(E\) be a real Euclidean space endowed with an orthogonal
complex structure \(J_E\), that is,
\[
 J_E^2=-I,
 \qquad
 J_E^{\!*}=-J_E.
\]
Define
\[
 \omega_J=\exp\!\left(\frac{2\pi}{3}J_E\right),
 \qquad
 \omega_J^3=I,
 \qquad
 I+\omega_J+\omega_J^2=0.
\]
On \(E\otimes\mathbb R^3\), the internal triadic transform is
\begin{equation}
 \mathcal F_3^{(J)}
 =\frac1{\sqrt3}\sum_{r,s=0}^{2}
   \omega_J^{rs}\otimes|r\rangle\langle s|.
 \label{eq:fourier-triadica-interna}
\end{equation}

\begin{proposition}[Orthogonality of the internal triadic transform]
\label{prop:fourier-triadica-unitaria}
The map \(\mathcal F_3^{(J)}\) is orthogonal, and its inverse is obtained
by replacing \(J_E\) with \(-J_E\):
\[
 \bigl(\mathcal F_3^{(J)}\bigr)^{-1}
 =\mathcal F_3^{(-J)}.
\]
In a chart that identifies \(J_E\) with multiplication by \(i\), it recovers the
complex Fourier matrix of order three.
\end{proposition}

\begin{proof}
The orthogonality of \(J_E\) implies
\((\omega_J^k)^{\!*}=\omega_J^{-k}\). The \((r,r')\) block of
\(\mathcal F_3^{(J)}\mathcal F_3^{(-J)}\) is
\[
 \frac13\sum_{s=0}^{2}\omega_J^{(r-r')s}.
\]
The sum equals \(I\) if \(r=r'\) and equals zero in the other two cases by
\(I+\omega_J+\omega_J^2=0\). The composition is therefore the identity.
The final assertion follows from
\(\exp(\theta J_E)\leftrightarrow e^{i\theta}\) in the complex chart.
\end{proof}

This construction makes the elliptic type of TRIT operational in Bell
instruments. A local setting may be composed from
\(\mathcal F_3^{(J)}\), diagonal phase transports, and projectors onto the ternary
basis. The orientation and memory transports of each wing must determine
those phases before the inequality is evaluated. Thus, \(J_E\) supplies the internal
rotation, whereas the genealogical register supplies the genealogical selection of the instruments.

\subsection{Joint Representation and Independence of the Marginals}

The confluence with quantum mechanics is expressed through a realization map.
\begin{equation}
 \mathcal R_{\rm Bell}:X_{\rm enr}\longrightarrow
 \mathcal D(\mathcal H_A\otimes\mathcal H_B),
 \label{eq:realizacion-bell}
\end{equation}
where \(\mathcal D(\mathcal H_A\otimes\mathcal H_B)\) is the space of
density operators. Let
\(\{M_{A,k}^{a}\}_{k\in\mathbb Z/3\mathbb Z}\) and
\(\{N_{B,\ell}^{b}\}_{\ell\in\mathbb Z/3\mathbb Z}\) be local POVM operators. The
joint rule is
\begin{equation}
 P(k,\ell\mid a,b)
 =\int_{X_{\rm enr}}
  \operatorname{tr}\!\left[
   \mathcal R_{\rm Bell}(x)
   \bigl(M_{A,k}^{a}\otimes N_{B,\ell}^{b}\bigr)
  \right]\dd\mu(x).
 \label{eq:born-bell-hmt}
\end{equation}

\begin{proposition}[Local Completeness Implies Independence of the Marginals]
\label{prop:no-signaling-hmt}
If
\[
 \sum_{\ell=0}^{2}N_{B,\ell}^{b}=I_B
 \quad\text{for every }b,
\]
then Alice's marginal is independent of Bob's setting. The symmetric
statement holds for the other wing.
\end{proposition}

\begin{proof}
Upon summing \eqref{eq:born-bell-hmt} over \(\ell\), completeness replaces
\(\sum_\ell N_{B,\ell}^{b}\) by \(I_B\). One obtains
\[
 \sum_{\ell}P(k,\ell\mid a,b)
 =\int\operatorname{tr}\!\left[
   \mathcal R_{\rm Bell}(x)
   \bigl(M_{A,k}^{a}\otimes I_B\bigr)
  \right]\dd\mu(x),
\]
an expression independent of \(b\).
\end{proof}

Equation \eqref{eq:realizacion-bell} fixes the type of the intertwiner between the
HMT genealogy and the bipartite space of the physical realization. When
\(\mathcal R_{\rm Bell}(x)\) is separable for almost every \(x\), the distribution
can be rewritten through shared variables and remains in the local domain.
A nonseparable realization may leave that factorization and preserve, by the
preceding proposition, the independence of the marginals from the remote setting.
Global correlation and signal transmission are therefore distinct properties of the
same construction.

\subsection{Holographic Reconstruction of Correlation via Indistinguishable Path Classes}

The HMT structure now organizes the local--global question through typed objects.
The whole is a compatible section of the extension system,
with memory of its genealogy. Each experimental wing is a boundary
equipped with an instrument that publishes a fiber. The joint preparation fixes the
common state; the instruments fix the local questions; the Born rule fixes
the probabilities; CHSH or CGLMP determines whether those probabilities admit a
local decomposition.

Holographic correlation is realized as a reading of the same global state within
a joint representation. Independence of the marginals ensures that each
local statistic remains invariant under the remote setting, while
nonseparability may
prevent the complete distribution from being reduced to independent responses
conditioned on a classical variable.

A precise probative hierarchy is thereby established. The local CHSH and CGLMP bounds, persistence of the bound under enrichment of the state, invariance under change to \(\log 3\), the transform \(\mathcal F_3^{(J)}\), and independence of the marginals by completeness of the POVM are closed mathematical results. The HMT architecture of Bell is fixed by
\[
 X_{\rm enr}
 \xrightarrow{\ \mathcal R_{\rm Bell}\ }
 \mathcal D(\mathcal H_A\otimes\mathcal H_B)
 \xrightarrow{\ M^a\otimes N^b\ }
 P(k,\ell\mid a,b)
 \xrightarrow{\ \mathrm{CGLMP}\ }
 I_3^{\rm CGLMP}.
\]
The evaluation of a concrete HMT violation is obtained by constructing
\(\mathcal R_{\rm Bell}\) and the four instruments from the enriched genealogical record,
certifying their independence from adjustments, and calculating the distribution without
postselection. The chapter provides the domain, the operators, the bounds, and the
realization criterion; the numerical evaluation remains reserved for the distribution
derived from those instruments.
